% Laget automatisk av elevark/lag_elevark.py fra sim/magnet-og-spole.html. Endringer her blir overskrevet.
\documentclass[11pt]{article}
\usepackage{elevark}
\begin{document}
\elevarkhode{Fysikk 2 · Elektromagnetisk induksjon}{Magnet og spole}{https://fysikklab.no/sim/magnet-og-spole.html}{Simuleringen}

Dra magneten inn og ut av spolen, eller la den kjøre gjennom med jevn fart. Når fluksen gjennom spolen endrer seg, blir det indusert en spenning. Hvilken vei går strømmen, og hvorfor blir det to pulser med motsatt fortegn?

\section*{Oppgaver}

\begin{oppgaveliste}
\oppgave \fase{gjett}{Gjett}Magneten ligger i ro midt inne i spolen. Blir det indusert en spenning? Hva om den ligger i ro rett utenfor?
\svar
\oppgave \fase{test}{Test}Dra magneten sakte inn i spolen, og så fort. Hva skjer med spenningen? Hold den i ro inne i spolen. Hva viser voltmeteret?
\svar
\oppgave \fase{forklar}{Forklar}Trykk Start. Forklar hvorfor spenningen har motsatt fortegn når magneten er på vei inn og på vei ut. Hvorfor er $ε = 0$ akkurat når magneten er midt i spolen?
\svar
\oppgave \fase{forklar}{Forklar}Se på polene til spolen når magneten nærmer seg og når den fjerner seg. Bruk Lenz' lov og forklar hvorfor spolen alltid motvirker bevegelsen.
\svar
\oppgave \fase{test}{Undersøk}Hvordan avhenger toppverdien av $ε$ av farten og av antall vindinger? Hva skjer med arealet under hver puls når farten dobles?
\svar
\oppgave \fase{test}{Regn}Med $N = 200$ og sterk magnet er største fluks $NΦ = 40$ mWb. Med farten 1,0~m/s bruker magneten 19~ms fra fluksen er halvparten av den største til den er størst. Regn ut gjennomsnittsspenningen i dette tidsrommet, og sammenlign med toppverdien.
\svar
\oppgave \fase{utfordring}{Utfordring}Spolen er koblet til en lyspære som lyser når magneten går gjennom. Hvor kommer energien fra? Hvorfor ville det vært umulig om strømmen gikk motsatt vei?
\svar
\end{oppgaveliste}

\end{document}
